\section{Secret Sharing}

Consider: $s$ is the signature for a peace treaty.
We need three parties -- a general, a
judge, and the president -- to all agree
to sign.
Have each party receive shares
$s_1$, $s_2$, $s_3$ of the secret $s$.
Having one or two of the shares is not
enough to get $s$
If all three parties come together, they
can obtain $s$.

We have a trusted secret broker compute 
$s_1$ and $s_2$ randomly, and choose 
$s_3$ s.t. $s_1 \oplus s_2 \oplus s_3 = s$. 
Out of 3 shares, you need 3 to recover the secret.
This can be generalized to $n$-out-of-$n$ secret 
sharing. 

In 1979, Adi Shamir introduced a
$t$-out-of-$n$ secret sharing scheme where
any $t$ shares (but no fewer) could 
reconstruct the shared secret. 
Shamir reasoned that $t$ points 
are required to specify a $t - 1$-degree 
polynomial. However, with $t-1$ points, 
there are infinitely many polynomials 
which fit the points. If the secret is, 
for instance, the coefficients of the 
polynomial, or the value of the $y$ 
intercept. For the latter, it suffices 
to sample the polynomial at random 
points to generate each of the secret 
shares. 

One $t$ shares are obtained, we can 
use Lagrange interpolation to 
recover the polynomial. 
Any subset of up to $t - 1$ shares does not leak 
any information about the secret.
This scheme is perfectly secure and 
very flexible! For instance, 
One can increase $n$ and add new shares 
without affecting other shares.
Existing shares can be removed without 
affecting other shares.
It is possible to replace all shares 
(or even change $t$) without changing the secret
and without revealing any information on the 
secret by selecting a new polynomial
$\hat{P}(x)$ and a new set of shares.
Some parties can be given more than one share.